\subsection{Semantic Prompt Effects, Ordering, and Attribution}
\label{app:prompt-response}

This secondary analysis was designed after inspecting the results. It uses paired generic and
pedagogy responses already rated for six historical models: MiniMax-M3,
MiniMax-M2.7, GLM-5.2, DeepSeek-V4-Pro, Doubao-Seed-2.0-Pro, and Qwen3.5-4B.
There are 79 standard and 40 hard MathDial contexts, each with both prompt
conditions and all six models. The eight semantic dimensions yield 11,424
response--dimension rows, not 11,424 independent replies. Analyses use the
saved three-annotator consensus. They do not revise the recorded decisions about the dimensions
or constitute independent prospective testing.

\paragraph{Prompt shift relative to default spread.}
For each task and dimension, the absolute mean paired pedagogy-minus-generic
shift is divided by the full range of model means under the generic prompt.
Table~\ref{tab:prompt-scale} reports the three dimensions that passed the
measurement reliability criterion, with 5,000 context-bootstrap intervals. These are
ratios of descriptive point estimates, not tests that prompts universally
matter more than model choice. A ratio near one means that the mean shift
under this instruction is comparable to the observed baseline range.

\begin{table}[t]
\centering\small
\begin{tabular}{llr}
\toprule
Task & Dimension & Shift / default range [95\% CI]\\
\midrule
Standard & Response complexity & 1.111 [0.785, 1.401]\\
Standard & Semantic elicitation & 1.148 [0.946, 1.437]\\
Standard & Assistance directness & 1.000 [0.830, 1.204]\\
Hard & Response complexity & 0.848 [0.585, 1.214]\\
Hard & Semantic elicitation & 0.987 [0.777, 1.301]\\
Hard & Assistance directness & 1.011 [0.758, 1.361]\\
\bottomrule
\end{tabular}
\caption{Prompt shifts relative to default model spread for the rated behavioral dimensions. CI denotes confidence interval. Intervals are descriptive context-bootstrap intervals, not simultaneous tests across dimensions.}
\label{tab:prompt-scale}
\end{table}

\paragraph{Response heterogeneity and scale boundaries.}
Let $s$ and $p$ denote the generic and pedagogy scores. For a pooled increase,
the adjusted change is $(p-s)/(5-s)$; for a pooled decrease, it is
$(s-p)/(s-1)$. Only contexts
where every model has nonzero room to change in the tested direction are retained. The retained
standard/hard counts are 41/19 for elicitation, 58/31 for directness, and 36/16
for complexity. Model labels are permuted within paired contexts 5,000 times;
Benjamini--Hochberg correction covers eight dimensions within each task.
Adjusted $q$ values for elicitation are 0.0003/0.0016 and for directness
0.0003/0.0155 (standard/hard); complexity is 0.0003/0.7023. Thus only
elicitation and directness among the reliably measured dimensions retain
heterogeneity in both variants after adjustment. The exclusion of boundary
contexts changes the analysis subset, and normalization addresses one
floor/ceiling explanation rather than every possible source of heterogeneity.

\paragraph{Changes in model ordering.}
Table~\ref{tab:prompt-order-example} illustrates the Doubao/GLM directness
comparison. Across the six models, directness ordering reverses for 5/15
comparable pairs in standard and 6/14 in hard contexts; pairs tied in either
condition are excluded. These counts describe sample-mean crossings, not separate
significant pairwise reversals. The generic-to-pedagogy Spearman correlations
are 0.371 and 0.203; their exact two-sided permutation $p$ values, 0.497 and
0.711, do not establish either preserved ordering or its absence.

\begin{table}[t]
\centering\small
\begin{tabular}{llrrr}
\toprule
Task & Model & Generic & Pedagogy & Paired change [95\% CI]\\
\midrule
Standard & Doubao-Seed-2.0-Pro & 3.90 & 1.58 & -2.32 [-2.62, -2.01]\\
Standard & GLM-5.2 & 2.97 & 1.87 & -1.10 [-1.38, -0.84]\\
Hard & Doubao-Seed-2.0-Pro & 3.98 & 2.10 & -1.88 [-2.30, -1.43]\\
Hard & GLM-5.2 & 3.08 & 2.48 & -0.60 [-1.00, -0.17]\\
\bottomrule
\end{tabular}
\caption{An illustrative change in directness ordering. Higher scores indicate more direct assistance. Intervals describe each model's paired prompt shift, not the difference between models or a test of rank reversal.}
\label{tab:prompt-order-example}
\end{table}

\paragraph{Residual model attribution across instructions.}
A standardized multinomial logistic classifier predicts the six model
identities from eight item-centered semantic scores, training under one
instruction and testing under the other. All eight dimensions are included;
classification success does not validate them individually as traits.
Pooled generic-to-pedagogy accuracy is 0.301 [0.264, 0.335]; the reverse is
0.297 [0.266, 0.326], compared with chance 0.167. The four cross-task,
cross-prompt directions range from 0.263 to 0.300, with lower interval bounds
above chance. Intervals use 2,000 test-context bootstrap samples conditional
on the fitted classifier. Pooled transfers reuse paired contexts across conditions;
cross-task transfers compare the two benchmark variants rather than a new
prospectively sampled source population. This is supplementary evidence of
joint behavioral discrimination, distinct from the action-probability
predictions fixed before response collection in Section~\ref{sec:stability}. It does not establish immutable
model identity or an internal personality mechanism.
